% =====================================================================
% Segunda parte – 1. Introducción y plan de pruebas
% =====================================================================
\section{Introducción}

Esta parte documenta las pruebas de ShiftCash. Seguimos a \textcite{myers2011} y \textcite{pressman2010} y organizamos las pruebas en cuatro niveles: unitarias, de integración, de sistema y de aceptación. El análisis de caja negra y caja blanca se concentra en las tres funcionalidades centrales: registrar transacción (F1, CU01), gestionar presupuesto (F2, CU02) y gestionar meta de ahorro (F3, CU03). Las pruebas automatizadas cubren además el resto de los requisitos.

Las pruebas se ejecutaron dos veces. La \textbf{primera ejecución} (25 de octubre de 2026, sobre el commit \codigo{0efe28b}) encontró los defectos D-01 a D-10 (Apéndice \ref{ap:defectos}). La \textbf{segunda ejecución} (21 de noviembre de 2026, rama \codigo{fix/defectos}) se hizo después de corregirlos y encontró además D-11. Cada sección muestra el resultado de ambas ejecuciones.

\section{Plan de pruebas}

\subsection{Objetivos}

\paragraph{Objetivo general.} Verificar que ShiftCash cumple sus requisitos funcionales y no funcionales, y encontrar y corregir los defectos que afectan la exactitud de los cálculos financieros, la seguridad y la experiencia del usuario.

\paragraph{Objetivos específicos.}
\begin{enumerate}[label=OP\arabic*., leftmargin=1.6cm]
  \item Verificar que los esquemas Pydantic rechazan montos, tipos, períodos y fechas inválidos con 422.
  \item Cubrir todos los caminos independientes de \codigo{_toca_hoy}, del cálculo del período mensual de un presupuesto y de \codigo{generate_daily_amounts}.
  \item Verificar la integración incremental de modelos, servicios, rutas y cliente en el flujo de cada caso de uso.
  \item Medir la seguridad, el rendimiento, la usabilidad, la portabilidad, la fiabilidad y la mantenibilidad frente a los RNF de la Tabla \ref{tab:rnf}.
  \item Validar con un usuario final, en la aplicación desplegada, que el sistema satisface los requisitos (porcentaje de RF aceptados).
  \item Alcanzar una cobertura de líneas del backend de al menos 80\,\% (RNF09).
\end{enumerate}

\subsection{Alcance}

\begin{table}[H]
\caption{Alcance del plan de pruebas}
\label{tab:alcance-pruebas}
\footnotesize
\begin{tabularx}{\linewidth}{L{2.8cm} L{5.4cm} Y}
\toprule
Módulo & Componentes bajo prueba & Objetivo \\
\midrule
F1 Transacciones & \codigo{TransactionCreate}, \codigo{routes.transactions}, \codigo{category_service}, \codigo{scheduler._toca_hoy}, \codigo{registrar_recurrentes}, \codigo{useTransaccionForm} & Validación de monto, tipo y categoría; recurrencias sin duplicados. \\
F2 Presupuestos & \codigo{BudgetCreate}, \codigo{budget_service} (\codigo{get_date_range}, \codigo{calculate_spent}, \codigo{percentage_used}), \codigo{routes.budgets}, \codigo{PresupuestoModal} & Rango por período, gasto, porcentaje real y estado «superado». \\
F3 Metas & \codigo{SavingGoalCreate}, \codigo{saving_service} (\codigo{generate_daily_amounts}, \codigo{calculate_progress}, \codigo{toggle_day}), \codigo{routes.savings}, \codigo{useMetaForm} & Exactitud del plan, marcado de días y progreso. \\
Soporte & \codigo{routes.users}, \codigo{routes.profile}, \codigo{routes.categories}, \codigo{auth} & Registro, inicio de sesión (con Google simulado), perfil, categorías y eliminación en cascada. \\
Transversal & \codigo{get_current_user}, \codigo{PrivateRoute}, CORS, despliegue & Seguridad, rendimiento, portabilidad, usabilidad, fiabilidad y mantenibilidad. \\
\bottomrule
\end{tabularx}
\tablenote{Fuera del alcance: la recuperación de contraseña y el inicio con Facebook, que no están implementados (Tabla \ref{tab:alcance}). Fuente: elaboración propia.}
\end{table}

\subsection{Estrategia: niveles, tipos, técnicas y herramientas}

\begin{table}[H]
\caption{Niveles, tipos, técnicas y herramientas de prueba}
\label{tab:estrategia}
\footnotesize
\begin{tabularx}{\linewidth}{L{1.9cm} L{2.5cm} L{3.8cm} Y}
\toprule
Nivel & Tipo & Técnica & Herramienta y justificación \\
\midrule
Unitario & Funcional y estructural & Clases de equivalencia, valores límite, camino básico \parencite{mccabe1976} & pytest \parencite{pytest2025} con pytest-cov: el backend es Python y la lógica crítica son funciones puras. \\
Integración & Incremental ascendente; basada en hilos & Pruebas de API con base de datos real (SQLite con FK) y flujo E2E & pytest + \codigo{TestClient} de FastAPI; Cypress \parencite{cypress2025} contra la API y la SPA en ejecución. \\
Sistema & Seguridad & Escaneo pasivo y activo (OWASP Top 10) & OWASP ZAP \parencite{zap2025,owasp2021}. \\
Sistema & Rendimiento y fiabilidad & Carga: 50 usuarios, 60 s de rampa, 10 iteraciones & Apache JMeter \parencite{jmeter2025}. \\
Sistema & Portabilidad & Compatibilidad entre navegadores y anchos & LambdaTest \parencite{lambdatest2025} y emulación móvil. \\
Sistema & Usabilidad & Inspección heurística y tarea cronometrada & Navegador y usuarios reales (PU-06). \\
Sistema & Mantenibilidad & Métricas estáticas & \codigo{docs/scripts/metricas.py}, ESLint, SonarQube \parencite{sonarqube2025}. \\
Aceptación & Funcional de usuario & Escenarios derivados de los RF, ejecutados por un usuario final & Guion y acta (Apéndice \ref{ap:acta}); automatización de apoyo con Selenium \parencite{selenium2025}. \\
\bottomrule
\end{tabularx}
\tablenote{Fuente: elaboración propia.}
\end{table}

\subsection{Ambiente de pruebas}

\begin{table}[H]
\caption{Ambiente de pruebas y versiones exactas}
\label{tab:ambiente}
\footnotesize
\begin{tabularx}{\linewidth}{L{3.6cm} Y}
\toprule
Elemento & Especificación \\
\midrule
Código bajo prueba & Repositorio \codigo{Amaya0910/ShiftCash}: commit \codigo{0efe28b} (primera ejecución) y rama \codigo{fix/defectos} (segunda ejecución). \\
Ambiente desplegado & Frontend en Vercel (\url{https://shifts-chash-version-amaya.vercel.app/}); API en Render; PostgreSQL en Supabase. \\
Ambiente local & Windows 11; API con Uvicorn 0.49 en \codigo{localhost:8000} y SQLite; frontend con Vite 8.0.9 en \codigo{127.0.0.1:5173}; Node.js 25.9.0. Datos ficticios: un usuario de prueba, 4 categorías, 6 transacciones de noviembre de 2026, 2 presupuestos y 2 metas. \\
Unitarias e integración & Python 3.12.10, pytest 9.2.0, pytest-cov 7.1.0 (coverage 7.17.4), httpx 0.28.1; SQLite en memoria con claves foráneas activadas. \\
Integración E2E & Cypress 13.17.0 (Electron 27.3.10, sin interfaz). \\
Aceptación automatizada & Selenium 4.53.0 (Python) con Google Chrome 157.0.8103.80, ventana de 1366 × 900 px. \\
Rendimiento & Apache JMeter 5.6.3 sobre Eclipse Temurin JDK 21.0.13. \\
Seguridad y calidad & OWASP ZAP 2.16.1 (Docker); SonarQube Community 26.7 (primera ejecución). \\
Portabilidad & LambdaTest Real Time: Chrome 154 en macOS y Firefox 155 en Windows 11. \\
\bottomrule
\end{tabularx}
\tablenote{Fuente: elaboración propia.}
\end{table}

\subsection{Recursos asignados}

\begin{table}[H]
\caption{Recursos humanos, técnicos y de tiempo}
\label{tab:recursos}
\footnotesize
\begin{tabularx}{\linewidth}{L{3.4cm} L{4.2cm} Y}
\toprule
Recurso & Rol o uso & Asignación \\
\midrule
Isabella Amaya & Diseño y ejecución de pruebas & Todos los tipos de prueba, en conjunto con Brayan \\
Brayan Vásquez & Diseño y ejecución de pruebas & Todos los tipos de prueba, en conjunto con Isabella \\
Usuario final & Ejecuta la prueba de aceptación y firma el acta & \resUsuarioFinal \\
5 participantes & Tarea cronometrada PU-06 & \resPUParticipantes \\
Equipo de cómputo & Ejecución local y herramientas & Computador del equipo con Windows 11 \\
Herramientas & Tabla \ref{tab:ambiente} & Todas gratuitas o de código abierto (LambdaTest en plan gratuito) \\
Tiempo & Primera ejecución & 25 de octubre de 2026 \\
Tiempo & Corrección y segunda ejecución & 20 y 21 de noviembre de 2026 \\
Tiempo & Aceptación con usuario final & \resFechaSesion \\
Horas dedicadas & Total del plan de pruebas & Estimadas: 15 h por semana por autor entre el 13 de agosto y el 21 de noviembre de 2026 (unas 14 semanas), cerca de 420 horas \\
\bottomrule
\end{tabularx}
\tablenote{Fuente: elaboración propia.}
\end{table}

\subsection{Criterios de aprobación}

\begin{itemize}
  \item \textbf{Unitarias e integración:} el 100\,\% de las pruebas automatizadas pasan y la cobertura de líneas del backend es de al menos 80\,\%.
  \item \textbf{Sistema:} cada RNF alcanza su valor objetivo (Tabla \ref{tab:rnf}); ninguna alerta de ZAP de riesgo alto.
  \item \textbf{Aceptación:} al menos el 90\,\% de los RF aceptados por el usuario final y ningún defecto abierto de severidad crítica o alta.
  \item \textbf{Suspensión:} se suspende la ejecución si la API no está disponible; se reanuda al restablecerla.
\end{itemize}
